% ============================================================
%  CSE200 Online-1 (A2)  ->  "Mathematical Foundations of Image Synthesis"
%  Compile:  pdflatex Q1_ImageSynthesis.tex
%  Needs the `mwe` package for the example-image-a/b/c demo images.
% ============================================================
\documentclass[11pt]{article}

\usepackage[T1]{fontenc}
\usepackage[utf8]{inputenc}
\usepackage[margin=1in]{geometry}
\usepackage{amsmath}      % equation, align, cases, split
\usepackage{amssymb}
\usepackage{graphicx}
\usepackage{subcaption}   % subfigures
\usepackage{multirow}     % multirow table cells
\usepackage{mwe}          % provides example-image-a / -b / -c

\title{Mathematical Foundations of Image Synthesis}
\author{Your Name (Your Student ID)}
\date{January 4, 2026}

\begin{document}
\maketitle

% ------------------------------------------------------------
\section{Introduction}
Image synthesis is the process of generating images from mathematical and
algorithmic descriptions. Modern graphics systems rely on \textbf{precise
mathematical models} and efficient computation. \textit{Even small numerical
errors} can result in visible artifacts in rendered images.

Typical formulations involve scalar values such as $x$, indexed parameters like
$I_{ij}$, and exponential terms $e^{x^2}$. These elements are combined to
simulate realistic lighting behavior.

% ------------------------------------------------------------
\section{Core Components}

\subsection{Rendering Pipeline}
The rendering pipeline consists of the following stages:
\begin{itemize}
    \item Scene Definition
    \item Camera Configuration
    \begin{itemize}
        \item Light Placement
    \end{itemize}
    \item[\textbf{Important}] Shading Model
    \begin{itemize}
        \item Diffuse Shading
        \item Specular Highlights
    \end{itemize}
\end{itemize}

\subsection{Execution Order}
\begin{enumerate}
    \item Input Processing
    \item Ray Construction
    \begin{enumerate}
        \item Primary Rays
        \item Secondary Rays
    \end{enumerate}
    \item Intersection Evaluation
    \begin{itemize}
        \item Object Hit Test
    \end{itemize}
\end{enumerate}

% ------------------------------------------------------------
\section{Mathematical Modeling}
Light interaction with surfaces can be expressed using the following equations.
\begin{equation}
    L = I_{ij}\cdot \cos(\theta) + R^2
\end{equation}
\[
    k = \frac{ab}{cd}
\]
\begin{equation}
    \begin{split}
        f(x) &= x^2 + 2x + 1 \\
             &= (x+1)^2
    \end{split}
\end{equation}
\begin{align*}
    |y| &= \begin{cases} y  & \text{if } y \ge 0 \\ -y & \text{if } y < 0 \end{cases}\\[1ex]
    p(t) &= \begin{cases}
                0,            & t < 0, \\
                e^{-t},       & 0 \le t < 3, \\
                \ln(t-2),     & t \ge 3
            \end{cases}\\[1ex]
    \mathcal{N}(x_i,\mu,\sigma) &= \frac{1}{\sqrt{2\pi}\,\sigma}\,
        e^{-\frac{(x_i-\mu)^2}{2\sigma^2}}
\end{align*}

% ------------------------------------------------------------
\section{Performance Comparison}
\begin{table}[ht]
    \centering
    \begin{tabular}{|l|c|c|}
        \hline
        Method & Accuracy & Time Complexity \\
        \hline
        Rasterization & Medium & $O(n)$ \\
        \hline
        \multirow{2}{*}{Ray Tracing} & High      & $O(n^2)$ \\
                                     & Very High & $O(n^3)$ \\
        \hline
        \multicolumn{2}{|c|}{Hybrid Method} & $O(n^2)$ \\
        \hline
    \end{tabular}
    \caption{Comparison of Image Synthesis Techniques}
    \label{tab:perf}
\end{table}

The results in Table~\ref{tab:perf} highlight the trade-off between accuracy and
computational cost.

% ------------------------------------------------------------
\section{Visual Illustration}
\begin{figure}[ht]
    \centering
    \begin{subfigure}{0.3\textwidth}
        \centering
        \includegraphics[width=\textwidth]{example-image-a}
        \caption{Output A}
    \end{subfigure}
    \hfill
    \begin{subfigure}{0.3\textwidth}
        \centering
        \includegraphics[width=\textwidth]{example-image-b}
        \caption{Output B}
    \end{subfigure}
    \hfill
    \begin{subfigure}{0.3\textwidth}
        \centering
        \includegraphics[width=\textwidth]{example-image-c}
        \caption{Output C}
    \end{subfigure}
    \caption{Comparison of three rendering outputs.}
    \label{fig:outputs}
\end{figure}

Figure~\ref{fig:outputs} compares three different rendering outputs arranged in a
single row.

% ------------------------------------------------------------
\section{Conclusion}
Image synthesis combines \underline{mathematical precision}, algorithmic design,
and efficient implementation. As scene complexity increases, choosing the
appropriate method becomes critical for balancing quality and performance.

\end{document}
